\documentclass[11pt]{article}

\usepackage[margin=1in]{geometry}
\usepackage{amsmath}
\usepackage[T1]{fontenc}
\usepackage{lmodern}

\title{A Reproducible Demonstration of a Thresholded Risk Model}
\author{Demo Research Group}
\date{\today}

\begin{document}

\maketitle

\begin{abstract}
This short paper accompanies a local research workspace used to test
source-to-PDF selection, repository-aware editing, and approval-based LaTeX
changes.  We describe a deliberately small nonlinear risk model whose numerical
constants are shared with an executable Python implementation.  The example is
simple enough to inspect quickly while still containing equations, citations,
cross-references, a figure, and an included section.
\end{abstract}

\section{Introduction}
Reproducible writing is easiest when a manuscript and its computational sources
can be inspected together.  A reader may notice a numerical assumption in the
rendered paper, select the relevant sentence, and ask whether the implementation
uses the same value.  Conversely, a researcher who changes a constant in code
should be able to locate every sentence and equation that depends on it.

Our example follows the broad principle that assumptions should be explicit and
checked against the analysis that produced the reported result
\cite{gelman2013bayesian}.  It is not intended as a scientific claim.  Instead,
it provides a compact fixture for testing coordinated edits.  Section~\ref{sec:method}
defines the model, Figure~\ref{fig:pipeline} summarizes the workflow, and
Section~\ref{sec:discussion} identifies useful modification tasks.

\section{Method}
\label{sec:method}

For observation $i$, let $t_i$ denote a temperature anomaly in degrees Celsius
and let $r_i$ denote a standardized rainfall index.  The analysis first centers
temperature using the offset $\delta=0.8$ and forms the linear signal
\begin{equation}
  \eta_i = 1.7\,(t_i-\delta) + 0.6r_i
    + \log\!\left(\frac{0.12}{1-0.12}\right).
  \label{eq:linear}
\end{equation}
The constants $1.7$, $0.8$, and $0.12$ correspond respectively to
\texttt{SIGMOID\_\allowbreak STEEPNESS},
\texttt{TEMPERATURE\_\allowbreak OFFSET\_\allowbreak C}, and
\texttt{BASELINE\_\allowbreak RISK} in \texttt{analysis/model.py}.  Keeping
these values visible in prose makes cross-file consistency straightforward to
test.

We transform the signal with the logistic function
\begin{equation}
  p_i = \frac{1}{1+\exp(-\eta_i)},
  \label{eq:risk}
\end{equation}
so $0<p_i<1$.  An intervention is recommended when $p_i\geq 0.65$, matching
the code constant \texttt{DECISION\_\allowbreak THRESHOLD}.  The implementation
also records the seed value 17 for deterministic downstream examples, although
the score in Equation~\eqref{eq:risk} itself contains no random component.

\begin{figure}[htbp]
  \centering
  \fbox{\parbox{0.22\linewidth}{\centering
    \textbf{Inputs}\\[3pt]
    $t_i$ and $r_i$}}
  \hspace{0.03\linewidth}
  \rule[1.2ex]{0.08\linewidth}{0.5pt}
  \hspace{0.03\linewidth}
  \fbox{\parbox{0.22\linewidth}{\centering
    \textbf{Risk model}\\[3pt]
    $p_i=\operatorname{logit}^{-1}(\eta_i)$}}
  \hspace{0.03\linewidth}
  \rule[1.2ex]{0.08\linewidth}{0.5pt}
  \hspace{0.03\linewidth}
  \fbox{\parbox{0.22\linewidth}{\centering
    \textbf{Decision}\\[3pt]
    intervene if $p_i\geq0.65$}}
  \caption{The demo pipeline from observed inputs to a thresholded decision.
  The diagram uses only native LaTeX boxes and rules, so no generated image
  asset is required.}
  \label{fig:pipeline}
\end{figure}

This construction resembles a basic generalized linear model, but its role here
is pedagogical.  The parameter values were chosen to create several meaningful
selection targets rather than estimated from real observations.  A useful test
prompt is to select the threshold sentence and request a simultaneous change to
the manuscript and Python implementation.


\section{Discussion}
\label{sec:discussion}
The model is intentionally sensitive to small editorial changes.  For example,
changing the offset from $0.8\,^{\circ}\mathrm{C}$ to a different calibration
value requires an update to both the Python constant and Equation~\eqref{eq:linear}.
Changing the decision threshold likewise affects the implementation and the
interpretation of the final classification.  These dependencies make the
workspace suitable for testing whether an editing agent can reason beyond the
currently selected paragraph.

The example also supports PDF-first interaction.  Selecting the displayed
equation should map back to the corresponding source lines, while selecting the
caption of Figure~\ref{fig:pipeline} should identify an ordinary prose region.
After an accepted edit, recompilation can regenerate the PDF and SyncTeX index;
after a rejected edit, the source should remain unchanged.

\section{Conclusion}
This demo combines a multi-file LaTeX paper with a small analysis program.  Its
purpose is practical: it provides enough structure to exercise synchronized
selection, citation handling, repository context, proposed diffs, approval, and
undo without introducing external data or nonstandard build dependencies.

\bibliographystyle{plain}
\bibliography{references}

\end{document}
